\section{Introduction}
\label{sec:introduction}

\paragraph{Motivation.}
Computer-aided impossibility results can be made auditable: a finite atlas can
record the residual cases that were generated, solved, and checked.  Such
evidence is about truth and certificates.  A repair recommendation is a
different object.  It says which assumptions, axioms, or implementation levers
should be removed to restore feasibility, and it therefore depends on the
interface at which the deletion is named.

This distinction matters even when the impossibility evidence is trusted.
M1 established an auditable finite-atlas workflow separating Proof, Audit,
Witness, and Assumption obligations.  M1.5 then showed that clause-multiset
equivalence is not enough to make raw lever-level repairs canonical:
bundled and split realizations can expose different inclusion-minimal raw
deletions.  M3 asks the positive follow-up question: when may those raw
implementation repairs nevertheless be reported canonically at a declared
contract interface?

\paragraph{Thesis.}
A valid impossibility certificate does not determine a valid repair report.
Reportability is contract-relative: the paper must say which contract atoms
are reportable, how implementation levers map to those atoms, and why grouped
reports preserve the reference residual semantics.  M3 gives sufficient
conditions, and under a monotone reference contract an exactness condition, for
when \(\GroupedRepair\) coincides pointwise with \(\ContractRepair\).
Candidate B supplies the motivating artifact-defined instantiation.

\begin{figure}[t]
  \centering
  \fbox{\parbox{0.91\textwidth}{\centering
    \textbf{Figure placeholder: M1.5 to M3 spine}\\[0.5em]
    M1.5 raw non-canonicity
    $\longrightarrow$ M3-A / M3-B / M3-C
    $\longrightarrow$ Candidate B grouped canonicity}}
  \caption{The paper spine separates the negative raw-level result from the
  positive contract-level theory.}
  \label{fig:spine}
\end{figure}

\paragraph{Contributions and boundary.}
The paper does not claim novelty for the broad phenomenon that MUS, MCS, or
repair families can vary with representation.  That sensitivity is part of the
background diagnosis and repair-enumeration landscape.  The positive
contribution is to specify when such low-level variation may be soundly
reported through a higher-level contract interface:
\begin{enumerate}[label=(\roman*)]
  \item explicit reportability contracts;
  \item M3-A, an atomic syntactic sufficient condition;
  \item M3-B, a non-atomic \(\GroupSoundness\) condition;
  \item M3-C, exactness under deletion-monotone reference contracts;
  \item a Lean-checked abstract finite-set theorem core; and
  \item a Candidate B artifact-backed instantiation of M3-B.
\end{enumerate}

M3-B recovers report-level canonicity on Candidate B only after the
reportability contract is made explicit and \(\GroupSoundness\) is verified.
It does not erase M1.5, make raw repairs canonical, or make M3-A applicable to
the non-atomic Candidate B realization.

\paragraph{Claim discipline.}
Table~\ref{tab:claim-discipline} records the qualifications that govern the
manuscript.

\begin{table}[t]
\centering
\scriptsize
\begin{tabularx}{\textwidth}{@{}p{0.32\textwidth}p{0.08\textwidth}X@{}}
\toprule
Claim & Allowed? & Required qualification \\
\midrule
M1.5 shows raw repair non-canonicity & yes & Candidate B / declared artifact scope \\
M3-A recovers Candidate B & no & M3-A does not apply to non-atomic split \code{minlib} \\
M3-B recovers report-level canonicity on Candidate B & yes & only under artifact-defined contract and verified \(\GroupSoundness\) \\
Candidate B raw repairs become canonical & no & raw non-canonicity remains \\
\(\GroupSoundness\) characterizes grouped correctness & yes & only for residually faithful realizations over \(\Psi\)-deletion-monotone contracts \\
Atomicity is necessary & no & M3-A is a sufficient condition only \\
Raw transport holds under M3-B & no & raw transport is an M3-A result only \\
M3 results extend to family scale & no & Candidate B is artifact-defined / base-scope unless separately lifted \\
Arrow provides a second M3 example & no & Arrow is a prerequisite testbed unless per-residual transfer is established \\
Artifact-defined contract equals semantic/Lean contract & no & semantic/formal upgrading is a separate truth/certificate obligation \\
Encoding-granularity sensitivity of repairs is new by itself & no & M3's novelty is contract-relative reportability and exactness conditions \\
\bottomrule
\end{tabularx}
\caption{Claim discipline for the M3 paper.}
\label{tab:claim-discipline}
\end{table}

\paragraph{Paper organization.}
Section~\ref{sec:contracts} defines contracts and repair predicates.
Sections~\ref{sec:atomic}--\ref{sec:monotone} state M3-A, M3-B, and M3-C.
Section~\ref{sec:candidate-b} records the artifact-backed instantiation.
Sections~\ref{sec:related}--\ref{sec:discussion} position and delimit the result.
Exact Lean signatures and artifact audit tables appear in the appendices.
